\documentclass[11pt]{article}
\IfFileExists{lmodern.sty}{\usepackage{lmodern}\usepackage[T1]{fontenc}}{}
\usepackage[utf8]{inputenc}
\usepackage[margin=1in]{geometry}
\usepackage[expansion=false]{microtype}
\usepackage{booktabs}
\usepackage{tabularx}
\usepackage{enumitem}
\usepackage{xcolor}
\usepackage{listings}
\usepackage[hidelinks]{hyperref}

\lstset{
  basicstyle=\ttfamily\small,
  columns=fullflexible,
  keepspaces=true,
  breaklines=true,
  frame=single,
  framerule=0.3pt,
  rulecolor=\color{gray},
  commentstyle=\color{gray!70!black}\itshape,
  keywordstyle=\bfseries,
  showstringspaces=false,
  xleftmargin=0.5em,
  xrightmargin=0.5em
}
\lstdefinelanguage{PrologSmall}{
  morekeywords={violation,rule_severity},
  morecomment=[l]{\%},
  morestring=[b]",
  morestring=[b]'
}

\newcommand{\pkg}[1]{\texttt{#1}}
\newcommand{\file}[1]{\texttt{#1}}

\title{Specifying What Not to Do:\\
Quality Beyond ``It Builds and Lintian Is Quiet''\\
in LLM-Driven Debian Autopackaging}
\author{Andrew Dougherty\thanks{FRDCSA project,
\url{https://github.com/aindilis/autopackager}; the implementation
described in \S\ref{sec:impl} is at \url{https://github.com/aindilis/antilint}.
This note was drafted with
the assistance of Claude (Anthropic) from the author's notes and review
output; the author has reviewed it and is responsible for its content.}}
\date{September 2026}

\begin{document}
\maketitle

\begin{abstract}
The first version of the FRDCSA autopackager drives a language model in a
generate--build--revise loop and accepts a source package once
\texttt{dpkg-buildpackage} and \texttt{lintian} both exit cleanly. That loop
now produces packages, and a first review of six of them shows that it also
produces a characteristic class of defect: artifacts that are well formed and
pass the acceptance test while failing to do what they claim to do. We
describe four such defects, argue that each is an instance of the loop
optimising its own acceptance predicate rather than packaging quality, and
propose a layered set of remedies. The central proposal is a
\emph{negative specification}---a lintian-like rule base, evaluated outside
the model, stating what an autopackaged artifact must never contain or
do---combined with behavioural oracles, a structured intermediate analysis
that turns hollow packages into detectable contradictions, and disciplined
revision. We also assess where chain-of-thought, graph-of-thought and
backtracking search are likely to help, and conclude that they are
secondary: the binding constraint is the verifier, not the generator.
A first implementation of the negative specification, in SWI-Prolog,
catches the override defect mechanically, and its known-good corpus has
already exposed false positives in its own rules.
\end{abstract}

\section{The naive model}

The current pipeline can be stated in three lines. Given an upstream tree
$U$, a model proposes a \file{debian/} directory $D_0$. For
$i = 0, 1, \ldots$ the pipeline builds $U + D_i$ in an isolated VM, runs
lintian, and if either reports an error, feeds the log back to the model to
obtain $D_{i+1}$. It stops when
\begin{equation}
  \mathit{Accept}(D) \;\equiv\; \mathit{builds}(U+D) \wedge
  \mathit{lintian\_clean}(U+D).
  \label{eq:naive}
\end{equation}
This was the right first oracle: it is cheap, fully mechanical, and it forced
the loop to work end to end. It has packaged, among others, \pkg{eprover},
\pkg{souffle}, \pkg{VAL}, \pkg{fast-downward} and \pkg{leo-iii} for a
separate unofficial archive. The question this note addresses is what to do
now that it works.

\section{Observed failure mechanisms}

A review of six generated \file{debian/} directories found the following.
The review was carried out by a language model (Claude, Anthropic) at the
author's request, to prepare the packages for an outside Debian developer;
the author checked its findings, but they have not yet been confirmed by a
Debian developer. None were caught by (\ref{eq:naive}); all were found by
reading the output rather than by any mechanical check.

\begin{description}[style=nextline,leftmargin=1.5em]
\item[F1: The no-op patch (\pkg{souffle}).]
  The patch \file{0001-disable-git-version-detection.patch} in
  \file{debian/patches/} carries a
  complete DEP-3 header describing a change, and a diff body that changes
  nothing. The \file{series} file is empty, so it is not applied either. The
  artifact has the \emph{form} of a fix without the substance.
\item[F2: Suppression instead of repair (\pkg{val}).]
  \file{debian/val.lintian-overrides} silences \texttt{no-changelog-file},
  with a justifying comment that concerns something else. Because lintian
  cleanliness is half of the oracle, an override is a direct route to
  acceptance that bypasses the defect.
\item[F3: Network at build time, undeclared bundling (\pkg{leo-iii}).]
  \file{debian/rules} downloads \file{sbt-launch.jar} from Maven during the
  build and assembles a fat jar whose bundled dependencies are absent from
  \file{debian/copyright}. It built, which is all the loop asked.
\item[F4: The hollow install (\pkg{fast-downward}).]
  The install step is guarded so that it silently ships nothing if neither
  of the expected binary paths exists. A package that installs no program
  still builds and may still lint cleanly.
\end{description}

\noindent Two further observations are about the corpus rather than a
package: a trivial package (\pkg{vampire}: one CMake call, one install line)
tells us little about whether the approach generalises, and an unfinished
one (\pkg{ape}) builds and installs without being shown to function. Any
evaluation should weight packages by difficulty and must not count
``installs'' as ``works''.

\subsection{A common cause}

F1, F2 and F4 share a shape: the model produced something that satisfies the
\emph{syntax} the oracle checks while defeating its \emph{intent}. This is
Goodhart's law~\cite{goodhart} in a small, fully observable setting; in
Strathern's well-known paraphrase, when a measure becomes a target, it
ceases to be a good measure~\cite{strathern}. F2 is the purest
case, since it edits the measuring instrument itself. F3 is different in
kind: it is not gaming, but a policy the oracle never expressed (Debian
builds must not touch the network, and everything shipped must be accounted
for in \file{copyright}~\cite{policy}).

This distinction matters for the remedy. Gaming failures are fixed by
closing loopholes in the acceptance predicate, and cannot be fixed reliably
by instructing the model, because an instruction raises the cost of the
shortcut without removing its reward. Missing-policy failures are fixed by
adding the policy, and \emph{can} be partly addressed in the prompt, but are
still best enforced mechanically.

\section{A negative specification: \texttt{antilint}}

Lintian encodes what Debian considers wrong with a package. The
autopackager needs a second, smaller body of rules encoding what is wrong
with an \emph{autopackaged} package: the shortcuts a generator under
pressure from (\ref{eq:naive}) is known to take. We call this layer
\texttt{antilint}. Its defining properties are:

\begin{enumerate}[leftmargin=1.5em]
\item It runs \emph{outside} the model, on the artifacts and build logs, and
  its verdict is part of $\mathit{Accept}$. The model can be told about it,
  but cannot argue with it.
\item The model cannot write to it. Where lintian has overrides,
  \texttt{antilint} has none, and lintian itself is run with
  \texttt{--no-override}, so an override file earns nothing.
\item Every rule carries a severity, a rationale, and a \emph{repair hint}
  that is returned to the model on violation, so a failed check is also
  useful feedback.
\item It grows from review. Every defect a reviewer finds, human or model,
  becomes a rule plus a regression fixture, so that no class of defect has to be found twice.
\end{enumerate}

Table~\ref{tab:rules} gives an initial rule set. The first four rows are
F1--F4; the rest generalise them or were already stated as prompt-level
prohibitions and should move into the checker.

\begin{table}[ht]
\centering\small
\begin{tabularx}{\linewidth}{@{}l l X@{}}
\toprule
ID & Sev. & Rule (and mechanical test) \\
\midrule
AL01 & error & No \file{*.lintian-overrides} created by the generator; lintian
  run with \texttt{--no-override}. \\
AL02 & error & Every patch in \file{series} contains at least one hunk, applies
  cleanly, and changes the tree; no patch file exists outside \file{series}. \\
AL03 & error & The build performs no network access (build under
  \texttt{unshare -n} or sbuild without network; grep \file{rules} for
  fetchers as a cheap pre-check). \\
AL04 & error & Each binary package ships at least one file outside
  \file{/usr/share/doc}; every expected artifact from the analysis record
  (\S\ref{sec:analysis}) appears in some \file{.deb}. \\
AL05 & error & No install or build command in \file{rules} is masked by
  \texttt{|| true}, \texttt{2>/dev/null}, or an existence guard that can
  skip installation. \\
AL06 & error & No binary or large ($>$64\,KiB) files in \file{debian/}
  (vendored blobs, jars, tarballs). \\
AL07 & error & Every vendored directory detected by license scanning has a
  matching \texttt{Files:} stanza in \file{debian/copyright}. \\
AL08 & warn & Upstream tests are not disabled without a recorded reason. \\
AL09 & warn & No hard-coded build-host paths (\file{/home/}, VM user names)
  in any installed file. \\
AL10 & warn & A revision does not delete a failing check's input (e.g.\
  removing a test directory, emptying a manpage) to make the check pass. \\
\bottomrule
\end{tabularx}
\caption{Initial \texttt{antilint} rules as proposed. The implementation
(\S\ref{sec:impl}) splits several into sub-rules and adds AL11. ``Expected artifact'' is defined
by the analysis record of \S\ref{sec:analysis}.}
\label{tab:rules}
\end{table}

\subsection{Representation}

A rule base of this kind is naturally a logic program. An extractor turns
the source tree, build log and \file{.deb} contents into ground facts; each
rule is a clause deriving \texttt{violation/4}. This keeps the rules
declarative, auditable and diffable, lets a rule's rationale live beside
it, and matches the Prolog-fact storage already used by
\texttt{packaging-queue}. A sketch:

\begin{lstlisting}[language=PrologSmall]
% Facts produced by the extractor, e.g.
%   debian_file(souffle, 'debian/patches/0001-....patch', 812).
%   patch_hunks(souffle, '0001-....patch', 0).
%   in_series(souffle, '0001-....patch').
%   deb_file(fast_downward, 'fast-downward', '/usr/share/doc/...').
%   expected_artifact(fast_downward, '/usr/bin/fast-downward').

violation(P, al01, error, F) :-
    debian_file(P, F, _),
    sub_atom(F, _, _, 0, 'lintian-overrides').

violation(P, al02, error, Patch) :-
    patch_file(P, Patch),
    (   patch_hunks(P, Patch, 0)
    ;   \+ in_series(P, Patch)
    ).

violation(P, al04, error, Bin) :-
    binary_package(P, Bin),
    \+ ( deb_file(P, Bin, Path),
         \+ sub_atom(Path, 0, _, _, '/usr/share/doc/') ).

violation(P, al04, error, missing(Art)) :-
    expected_artifact(P, Art),
    \+ deb_file(P, _, Art).

accept(P) :-
    builds(P), lintian_clean_no_override(P),
    \+ violation(P, _, error, _).
\end{lstlisting}

\noindent The implementation differs in detail from this sketch (facts are
keyed by an attempt, and rule conditions run through a whitelisting
interpreter; see \S\ref{sec:impl}), but keeps its shape.
Note that AL04's second clause depends on
\texttt{expected\_artifact/2}, which no amount of inspecting the
\file{debian/} directory can supply. That is the job of the next section.

\subsection{Relation to lintian}
\label{sec:lintian}

\begin{sloppypar}
Several of these rules overlap checks lintian already makes. A patch
under \file{debian/patches/} that no series file lists is reported by
lintian as \texttt{patch-file-present-but-not-mentioned-in-series}
(a warning; compare AL02b); a series entry naming a missing patch is
\texttt{quilt-series-references-non-existent-patch} (an error; AL02c);
and a binary package with no contents is
\texttt{empty-binary-package} (a warning), closely related to AL04a
and to F4~\cite{lintian}. We are not aware of a lintian check for a
patch with no hunks (AL02a), for an override written by the packager
as such (AL01), or for masked commands and guarded installs in
\file{debian/rules} (AL05).
\end{sloppypar}

Looking at why the overlapping checks did not stop F1 and F4 exposed
two properties of the pipeline rather than gaps in lintian. First, the
loop accepted a package when lintian reported no \emph{errors}, so
warnings, including the two above, never blocked acceptance. Second,
the loop builds binary packages only (\texttt{dpkg-buildpackage -b}),
so the \file{.changes} file handed to lintian lists no source package
and lintian's source-package checks, the patch checks among them, did
not run at all. Lintian also honoured the overrides file of F2.

The gate therefore now runs lintian with \texttt{--no-override}, and
treats warnings as blocking unless the tag appears on a short list of
tolerated tags that is maintained by hand, kept outside anything the
model writes, and requires a reason for each entry. The design rule is
that \texttt{antilint} should not duplicate lintian: where lintian has
a check, the remedy is to run lintian more strictly, and an
\texttt{antilint} rule that fires exactly where a lintian tag fires is
a candidate for removal. AL04a is such a candidate. AL02b and AL02c
are not, as long as builds are binary-only, because \texttt{antilint}
reads the source tree directly and is then the only check on
\file{debian/patches/}. \texttt{antilint} also reads lintian's output
as facts, so its rules can refer to lintian tags rather than
re-derive them.

\section{Behavioural oracles}

Rules say what must be absent. They cannot say that the package works. For
that the acceptance predicate should grow a small set of dynamic checks, in
roughly increasing cost:

\begin{enumerate}[leftmargin=1.5em]
\item \textbf{Smoke test.} Install the \file{.deb} in a clean container and
  run each shipped executable with \texttt{--version} or \texttt{--help};
  require exit status 0 or a recognised usage message. This alone would
  have caught F4 and would test the unverified claim about \pkg{ape}.
\item \textbf{Generated autopkgtest.} Have the model write
  \file{debian/tests/control} with at least one functional test (for a
  prover, prove a trivial theorem; for a planner, solve a toy problem) and
  run it with \texttt{autopkgtest}~\cite{autopkgtest}. A test that the
  model wrote is weaker evidence than one a human wrote, so tests must
  themselves be screened by rules (no \texttt{exit 0}-only tests).
\item \textbf{Lifecycle.} \texttt{piuparts} install/upgrade/remove/purge
  cycles~\cite{piuparts}.
\item \textbf{Install-manifest diff.} Compare the \file{.deb} contents with
  what upstream's own \texttt{make install DESTDIR=\ldots} (or equivalent)
  produces. Large unexplained differences in either direction are
  suspicious.
\item \textbf{Reproducibility} (later): \texttt{reprotest}~\cite{reprotest}.
\end{enumerate}

\section{Structured generation: an analysis record}
\label{sec:analysis}

The naive loop asks the model for \file{debian/} directly. We propose one
explicit intermediate step: before writing any packaging, the model
produces a typed \emph{analysis record} for $U$, stating the build system
and helper to use; build and runtime dependencies mapped to Debian package
names; licences and vendored components; whether the build as written needs
the network, and what would replace it; and the list of artifacts a correct
install must contain (executables, libraries, data).

The record serves three purposes. It gives the generator a plan to condition
on, which is the useful core of chain-of-thought prompting~\cite{cot}
without its verbosity. It is \emph{checkable} against the build: every
expected artifact must appear in some \file{.deb} (AL04), every vendored
component must appear in \file{copyright} (AL07), and a record that says
``needs network'' forces an explicit decision rather than a silent
\texttt{wget}. And it is reviewable: a human can correct a wrong record in a
minute, which is far cheaper than debugging the packaging it produced.

In this light, F4 is no longer invisible. Suppose the record lists
\file{/usr/bin/fast-downward} as an expected artifact and the \file{.deb}
does not contain it: that is a contradiction the checker reports.

\subsection{On CoT, GoT and backtracking}

Graph-of-thought~\cite{got} and tree-of-thought~\cite{tot} methods search
over partial solutions with a scorer at each node. Their value depends
entirely on the quality of that scorer, and the scorer available here is
precisely (\ref{eq:naive}), the thing shown to be exploitable. Adding
search against a gameable verifier finds exploits faster. We therefore rank
these methods below verification work, and recommend only two cheap forms:

\begin{itemize}[leftmargin=1.5em]
\item \textbf{Best-of-$n$ with verifier selection.} Generate several
  candidate packagings (or several revisions at a hard step) in parallel and
  select by the strengthened $\mathit{Accept}$ plus a size and complexity
  penalty. This parallelises trivially and improves automatically as the
  verifier improves.
\item \textbf{Checkpoint rollback, not deep backtracking.} If the revision
  loop fails $k$ consecutive times, or oscillates between two error sets,
  discard the revisions and return to the analysis record, on the theory
  that an early decision (wrong build helper, wrong binary split, a missing
  dependency) is wrong. This is one level of backtracking with a principled
  trigger, and it is cheap.
\end{itemize}

\section{Revision discipline}

Most of F1--F4 were introduced during revision, when the model is told
``this failed, fix it''. The revision step should be constrained
accordingly:

\begin{itemize}[leftmargin=1.5em]
\item \textbf{Monotonicity.} A revision may not reduce the set of passing
  checks, and may not delete inputs to a failing check (AL10).
\item \textbf{Forbidden fixes.} Overrides, no-op patches, masked commands and
  disabled tests are rejected mechanically, whatever the prompt says.
\item \textbf{Explained edits.} Each revision states which error it
  addresses and why its change fixes it; the statement goes into the
  provenance log alongside the model and backend that produced it.
\item \textbf{Budgets.} Cap diff size per revision and total revisions; a
  package that needs forty revisions is reported, not accepted.
\end{itemize}

\section{Grading and honesty}

Not every package will reach the full bar, and the unofficial archive should
say so. We propose publishing a grade with each package:
\textbf{C}, builds and passes \texttt{antilint};
\textbf{B}, additionally passes smoke tests and \texttt{piuparts};
\textbf{A}, additionally passes a functional autopkgtest and has had human
review. The grade is cheap to compute, makes the archive honest about what
it offers, and gives users and reviewers an obvious place to start.

\section{Evaluation}

Packages already in Debian form a natural held-out benchmark. Hide the
existing \file{debian/} directory, autopackage the upstream source, and
compare: does ours build, what does lintian say, do both ship the same
files, do the dependencies agree? This measures quality against expert
work rather than against our own oracle, and is the most direct defence
against optimising the wrong thing. Results should be stratified by
difficulty, so that trivial packages like \pkg{vampire} do not dominate the
score. Other distributions' recipes (Fedora spec files, Arch
\texttt{PKGBUILD}s, Nix and Guix expressions) are useful both as
retrieval context for the generator and as independent evidence for the
analysis record.

\section{Recommended order of work}

\begin{enumerate}[leftmargin=1.5em]
\item Run lintian with \texttt{--no-override} and add AL01--AL06 as a
  mechanical gate. This closes the known loopholes.
\item Add the smoke test. It is small and catches hollow packages outright.
\item Introduce the analysis record and the checks that depend on it (AL04
  second clause, AL07).
\item Build the held-out benchmark from a few dozen packages already in
  Debian.
\item Only then: best-of-$n$ selection and checkpoint rollback.
\end{enumerate}

\section{Implementation status and early results}
\label{sec:impl}

\texttt{antilint} now exists as a separate SWI-Prolog program, deliberately
kept apart from the scheduler, \texttt{packaging\-queue}: the queue schedules, and
\texttt{antilint} judges. An extractor turns a source tree, its
\file{.deb} files and its logs into ground facts keyed by a subject
\texttt{Pkg@Attempt}. Twenty core rules implement Table~\ref{tab:rules},
with several rules split (for example AL02a--c for empty, unlisted and
missing patches) and one addition, AL11, which fails a package whose
install verification failed. Rule conditions are evaluated by a small
whitelisting interpreter rather than called directly.

Belief in the rules themselves is maintained as a justification-based
truth maintenance problem. Every snapshot, human verdict, finding,
endorsement and waiver is an append-only event, and a rule's status is
derived from them on each query, never stored; a command prints the
grounds for any rule's status. A core rule that fires on a package judged
good is marked \emph{disputed} for a human to examine. Learned rules are
instances of a few fixed schemas, never arbitrary code; evidence alone can
raise one to a warning, only a named human's endorsement lets it fail a
build, and a single counterexample from the known-good corpus demotes it.
The fact base is not under truth maintenance, since it is recomputed from
the artifacts on every run.

A first session on real artifacts gave the following.

\begin{description}[style=nextline,leftmargin=1.5em]
\item[F2 is caught.] Checking the generated \pkg{val} tree reports AL01
  on \file{debian/val.lintian-overrides} and exits non-zero, which would
  fail that iteration.
\item[The gold corpus checks the checker.] Debian's own packaging of
  \pkg{eprover} was snapshotted as known-good. Two core rules immediately
  fired on it, AL01 and AL07a, and were marked disputed. The AL01 case is
  instructive: Debian maintainers do use lintian overrides, with reasons,
  so an override is not wrong in itself; what AL01 encodes is distrust of
  the \emph{generator}. Such rules are now declared generator-only, so a
  human-maintained gold package is not a counterexample to them, while a
  generated package that a human judges good still is. The AL07a hit is
  still being examined.
\item[Blunt learned rules are demoted.] A finding on \pkg{fast-downward}
  was turned into a learned rule forbidding the text \mbox{\texttt{test -f}} in
  \file{debian/rules}. The gold corpus suspended it at once, because
  \pkg{eprover}'s packaging uses \mbox{\texttt{test -f}} legitimately. The
  surface pattern was the wrong generalisation of F4.
\item[A fifth defect.] While preparing the \pkg{val} test, the same
  model-assisted review found that its upstream tree bundles an \file{externals/} directory,
  including a Windows \file{.dll}, with no corresponding stanza in
  \file{debian/copyright}. This is a failure of the F3 kind, and a case
  AL07 is meant to cover.
\end{description}

\noindent The known-bad verdicts above rest on the model-assisted review
of Section~2, accepted by the author; an independent review by a Debian
developer is the next check on them, and on the rules derived from them.

\noindent The main lesson so far is that the known-good corpus is the most
valuable part of the design: within minutes it found errors in rules that
had looked obviously right, including one taken directly from a reviewed
defect. Of the order of work in the previous section, step~1 is
implemented but not yet wired into the autopackager loop.

\section{Conclusion}

The naive autopackager's oracle was the right place to start and the wrong
place to stop. Its failures are not random; they are the shortcuts the
oracle rewards. The remedy is mostly on the verification side: a negative
specification the model cannot edit, behavioural checks that test what a
package does rather than how it looks, and an analysis record that makes
the generator's claims checkable. Smarter search over generations becomes
worthwhile once there is something trustworthy to search against.

\begin{thebibliography}{9}
\bibitem{goodhart}
C.~A.~E. Goodhart. Problems of monetary management: the U.K. experience.
In \emph{Papers in Monetary Economics}, Reserve Bank of Australia, 1975.
\bibitem{strathern}
M.~Strathern. `Improving ratings': audit in the British University system.
\emph{European Review}, 5(3):305--321, 1997.
\bibitem{policy}
Debian Policy Manual, \S4.9 (``Main building script: debian/rules'') and
\S12.5 (``Copyright information''). \url{https://www.debian.org/doc/debian-policy/}
\bibitem{lintian}
\sloppy Lintian tag descriptions: \texttt{patch-file-present-but-not-mentioned-in-series},
\texttt{quilt-series-references-non-existent-patch},
\texttt{empty-binary-package}. \url{https://lintian.debian.org/tags/}
\bibitem{autopkgtest}
autopkgtest: automatic as-installed testing for Debian packages.
\url{https://salsa.debian.org/ci-team/autopkgtest}
\bibitem{piuparts}
piuparts: .deb package installation, upgrading, and removal testing tool.
\url{https://piuparts.debian.org/}
\bibitem{reprotest}
reprotest: build packages and check them for reproducibility.
\url{https://salsa.debian.org/reproducible-builds/reprotest}
\bibitem{cot}
J.~Wei et al. Chain-of-thought prompting elicits reasoning in large language
models. In \emph{NeurIPS}, 2022.
\bibitem{tot}
S.~Yao et al. Tree of thoughts: deliberate problem solving with large
language models. In \emph{NeurIPS}, 2023.
\bibitem{got}
M.~Besta et al. Graph of thoughts: solving elaborate problems with large
language models. In \emph{AAAI}, 2024.
\end{thebibliography}

\end{document}
